% GENERATED FILE. Edit canonical lessons, then run npm run generate:books.
% canonical-source-hash: fnv1a64:9cc565288c5b1094
% canonical-lessons: KA-C344-ayudha, KA-C344-banduku, KA-C344-ranaranga, KA-C344-ratha, KA-C344-kadadu

\chapter{Weapons of Old}
\label{ch:ka-weapons-of-old}
\hlchaptermodality{344}

\begin{chapteropening}
\textbf{By the end of this chapter:} \emph{I can say a weapon, a gun, a battlefield, a chariot and to do battle.}

Complete the last lesson of chapter 344: I can say a weapon, a gun, a battlefield, a chariot and to do battle.
\end{chapteropening}

\section[\texorpdfstring{\kn{ಆಯುಧ} (āyudha)}{ಆಯುಧ (āyudha)}]{\kn{ಆಯುಧ} (āyudha) — a weapon}

\label{lesson:KA-C344-ayudha}



\begin{quote}
Before the new one: say the Kannada for an episode, then the Kannada for a background.
\end{quote}

\subsection*{You'll want to know: \kn{ಆಯುಧ}}
\textbf{\kn{ಆಯುಧ}} — \emph{āyudha} — \textquotedblleft{}a weapon\textquotedblright{}.

\begin{quote}
\kn{ಸೈನಿಕರು} \kn{ತಮ್ಮ} \kn{ಆಯುಧಗಳನ್ನು} \kn{ಕೆಳಗೆ} \kn{ಇಟ್ಟರು}.
\end{quote}

\begin{quote}
\emph{sainikaru tamma āyudhagaḷannu keḷage iṭṭaru.}
\end{quote}

\begin{quote}
The soldiers laid down their weapons.
\end{quote}

Plural: \kn{ಆಯುಧಗಳು}, \emph{āyudhagaḷu}.

Also a tool for work: \kn{ರೈತನ} \kn{ಆಯುಧ}.

\begin{cousinweb}
\textbf{\kn{ಆಯುಧ}} (\emph{āyudha}) is borrowed from Sanskrit \emph{āyudha}, \textquotedblleft{}a weapon\textquotedblright{} (CDIAL 1291), from the root \emph{yudh}, \textquotedblleft{}to fight\textquotedblright{}, as in \kn{ಯುದ್ಧ}.
\end{cousinweb}

\begin{grammarlens}[title={What you've built}]
One more word for telling what happened.
\end{grammarlens}

\subsection*{Your turn}
\begin{itemize}
\raggedright
  \item \emph{Say it:} \emph{āyudha}
  \item \emph{Say it:} \emph{āyudha}, once more
  \item \emph{Say it:} \emph{sainikaru tamma āyudhagaḷannu keḷage iṭṭaru.} — The soldiers laid down their weapons.
  \item \emph{Recall:} say the Kannada for an episode, then the Kannada for a background, then say \emph{āyudha} again
\end{itemize}

\begin{tcolorbox}[breakable,colback=teal!4,colframe=teal!35!black,title={Before you move on}]
What does \kn{ಆಯುಧ} mean? (\textquotedblleft{}A weapon\textquotedblright{}.) Say it once more.
\end{tcolorbox}
\section[\texorpdfstring{\kn{ಬಂದೂಕು} (bandūku)}{ಬಂದೂಕು (bandūku)}]{\kn{ಬಂದೂಕು} (bandūku) — a gun; a rifle}

\label{lesson:KA-C344-banduku}



\begin{quote}
Before the new one: say the Kannada for a background, then the Kannada for a weapon.
\end{quote}

\subsection*{You'll want to know: \kn{ಬಂದೂಕು}}
\textbf{\kn{ಬಂದೂಕು}} — \emph{bandūku} — \textquotedblleft{}a gun; a rifle\textquotedblright{}.

\begin{quote}
\kn{ಕಾವಲುಗಾರನ} \kn{ಕೈಯಲ್ಲಿ} \kn{ಹಳೆಯ} \kn{ಬಂದೂಕು} \kn{ಇತ್ತು}.
\end{quote}

\begin{quote}
\emph{kāvalugārana kaiyalli haḷeya bandūku ittu.}
\end{quote}

\begin{quote}
The watchman had an old gun in his hand.
\end{quote}

Plural: \kn{ಬಂದೂಕುಗಳು}, \emph{bandūkugaḷu}.

\begin{cousinweb}
\textbf{\kn{ಬಂದೂಕು}} (\emph{bandūku}) is a loan from Arabic \emph{bunduq}, which came through Persian and Hindustani \emph{bandūq}, \textquotedblleft{}a musket, a gun\textquotedblright{}.
\end{cousinweb}

\begin{grammarlens}[title={What you've built}]
One more word for telling what happened.
\end{grammarlens}

\subsection*{Your turn}
\begin{itemize}
\raggedright
  \item \emph{Say it:} \emph{bandūku}
  \item \emph{Say it:} \emph{bandūku}, once more
  \item \emph{Say it:} \emph{kāvalugārana kaiyalli haḷeya bandūku ittu.} — The watchman had an old gun in his hand.
  \item \emph{Recall:} say the Kannada for a background, then the Kannada for a weapon, then say \emph{bandūku} again
\end{itemize}

\begin{tcolorbox}[breakable,colback=teal!4,colframe=teal!35!black,title={Before you move on}]
What does \kn{ಬಂದೂಕು} mean? (\textquotedblleft{}A gun; a rifle\textquotedblright{}.) Say it once more.
\end{tcolorbox}
\section[\texorpdfstring{\kn{ರಣರಂಗ} (raṇaraṅga)}{ರಣರಂಗ (raṇaraṅga)}]{\kn{ರಣರಂಗ} (raṇaraṅga) — a battlefield}

\label{lesson:KA-C344-ranaranga}



\begin{quote}
Before the new one: say the Kannada for a weapon, then the Kannada for a gun.
\end{quote}

\subsection*{You'll want to know: \kn{ರಣರಂಗ}}
\textbf{\kn{ರಣರಂಗ}} — \emph{raṇaraṅga} — \textquotedblleft{}a battlefield\textquotedblright{}.

\begin{quote}
\kn{ಯುದ್ಧದ} \kn{ನಂತರ} \kn{ರಣರಂಗ} \kn{ಮೌನವಾಗಿತ್ತು}.
\end{quote}

\begin{quote}
\emph{yuddhada nantara raṇaraṅga maunavāgittu.}
\end{quote}

\begin{quote}
After the war the battlefield was silent.
\end{quote}

\kn{ರಣ}, battle, and \kn{ರಂಗ}, a stage: the stage of battle.

\begin{cousinweb}
\textbf{\kn{ರಣರಂಗ}} (\emph{raṇaraṅga}) is borrowed from Sanskrit, a compound of \emph{raṇa}, \textquotedblleft{}battle, war\textquotedblright{}, and \emph{raṅga}, \textquotedblleft{}a stage, an arena\textquotedblright{} (Monier-Williams).
\end{cousinweb}

\begin{grammarlens}[title={What you've built}]
One more word for telling what happened.
\end{grammarlens}

\subsection*{Your turn}
\begin{itemize}
\raggedright
  \item \emph{Say it:} \emph{raṇaraṅga}
  \item \emph{Say it:} \emph{raṇaraṅga}, once more
  \item \emph{Say it:} \emph{yuddhada nantara raṇaraṅga maunavāgittu.} — After the war the battlefield was silent.
  \item \emph{Recall:} say the Kannada for a weapon, then the Kannada for a gun, then say \emph{raṇaraṅga} again
\end{itemize}

\begin{tcolorbox}[breakable,colback=teal!4,colframe=teal!35!black,title={Before you move on}]
What does \kn{ರಣರಂಗ} mean? (\textquotedblleft{}A battlefield\textquotedblright{}.) Say it once more.
\end{tcolorbox}
\section[\texorpdfstring{\kn{ರಥ} (ratha)}{ರಥ (ratha)}]{\kn{ರಥ} (ratha) — a chariot}

\label{lesson:KA-C344-ratha}



\begin{quote}
Before the new one: say the Kannada for a gun, then the Kannada for a battlefield.
\end{quote}

\subsection*{You'll want to know: \kn{ರಥ}}
\textbf{\kn{ರಥ}} — \emph{ratha} — \textquotedblleft{}a chariot\textquotedblright{}.

\begin{quote}
\kn{ರಾಜ} \kn{ಬಿಳಿ} \kn{ಕುದುರೆಗಳ} \kn{ರಥದಲ್ಲಿ} \kn{ಬಂದನು}.
\end{quote}

\begin{quote}
\emph{rāja biḷi kuduregaḷa rathadalli bandanu.}
\end{quote}

\begin{quote}
The king came in a chariot drawn by white horses.
\end{quote}

Plural: \kn{ರಥಗಳು}, \emph{rathagaḷu}.

At temple fairs the god is pulled through the streets in a great wooden \kn{ರಥ}.

\begin{cousinweb}
\textbf{\kn{ರಥ}} (\emph{ratha}) is borrowed from Sanskrit \emph{ratha}, \textquotedblleft{}a chariot, a cart\textquotedblright{} (CDIAL 10602).
\end{cousinweb}

\begin{grammarlens}[title={What you've built}]
One more word for telling what happened.
\end{grammarlens}

\subsection*{Your turn}
\begin{itemize}
\raggedright
  \item \emph{Say it:} \emph{ratha}
  \item \emph{Say it:} \emph{ratha}, once more
  \item \emph{Say it:} \emph{rāja biḷi kuduregaḷa rathadalli bandanu.} — The king came in a chariot drawn by white horses.
  \item \emph{Recall:} say the Kannada for a gun, then the Kannada for a battlefield, then say \emph{ratha} again
\end{itemize}

\begin{tcolorbox}[breakable,colback=teal!4,colframe=teal!35!black,title={Before you move on}]
What does \kn{ರಥ} mean? (\textquotedblleft{}A chariot\textquotedblright{}.) Say it once more.
\end{tcolorbox}
\section[\texorpdfstring{\kn{ಕಾದಾಡು} (kādāḍu)}{ಕಾದಾಡು (kādāḍu)}]{\kn{ಕಾದಾಡು} (kādāḍu) — to do battle; to clash}

\label{lesson:KA-C344-kadadu}



\begin{quote}
Before the new one: say the Kannada for a battlefield, then the Kannada for a chariot.
\end{quote}

\subsection*{You'll want to know: \kn{ಕಾದಾಡು}}
\textbf{\kn{ಕಾದಾಡು}} — \emph{kādāḍu} — \textquotedblleft{}to do battle; to clash\textquotedblright{}.

\begin{quote}
\kn{ಎರಡು} \kn{ಸೈನ್ಯಗಳು} \kn{ಇಡೀ} \kn{ದಿನ} \kn{ಕಾದಾಡಿದವು}.
\end{quote}

\begin{quote}
\emph{eraḍu sainyagaḷu iḍī dina kādāḍidavu.}
\end{quote}

\begin{quote}
The two armies did battle all day.
\end{quote}

Past (I …): \kn{ಕಾದಾಡಿದೆ}, \emph{kādāḍide}.

Who one fights with takes -\kn{ಜೊತೆ} or -\kn{ಒಡನೆ}. Animals \kn{ಕಾದಾಡು} too.

\begin{cousinweb}
\textbf{\kn{ಕಾದಾಡು}} (\emph{kādāḍu}) is a native Dravidian compound: \emph{kādu}, \textquotedblleft{}to war, to fight, to contend\textquotedblright{} (DEDR 1447), and \emph{āḍu}, \textquotedblleft{}to move, to act\textquotedblright{} (DEDR 347).
\end{cousinweb}

\begin{grammarlens}[title={What you've built}]
One more word for telling what happened.
\end{grammarlens}

\subsection*{Your turn}
\begin{itemize}
\raggedright
  \item \emph{Say it:} \emph{kādāḍu}
  \item \emph{Say it:} \emph{kādāḍu}, once more
  \item \emph{Say it:} \emph{eraḍu sainyagaḷu iḍī dina kādāḍidavu.} — The two armies did battle all day.
  \item \emph{Recall:} say the Kannada for a battlefield, then the Kannada for a chariot, then say \emph{kādāḍu} again
\end{itemize}

\begin{tcolorbox}[breakable,colback=teal!4,colframe=teal!35!black,title={Before you move on}]
What does \kn{ಕಾದಾಡು} mean? (\textquotedblleft{}To do battle; to clash\textquotedblright{}.) Say it once more.
\end{tcolorbox}
